\subsection{API Evolution, Breaking-Change Detection, and Behavioural Compatibility}
\label{subsec:rw7}

The discipline of managing how a software interface changes over time is the empirical backdrop against which this article's compositional change-propagation mechanism must be positioned. The governing convention is semantic versioning~\cite{preston2013semver}, which encodes a behavioural promise in the version string: a major increment signals an incompatible change, a minor increment signals backward-compatible additions, and a patch increment signals backward-compatible fixes. A substantial body of empirical work has examined whether practice honors this promise. Raemaekers et al.~\cite{raemaekers2017semver} analysed seven years of release history in Maven Central and found that roughly one third of all releases introduce at least one breaking change, frequently without the corresponding major increment, undermining the version string as a compatibility signal. Ochoa et al.~\cite{ochoa2022breaking} later replicated and refined this analysis over 119{,}879 library upgrades and 293{,}817 clients, reporting a more optimistic 83.4\% compliance rate and a rising trend over time, but crucially observing that most breaking changes touch code no client actually uses---only a small fraction of clients are affected in practice. Decan et al.~\cite{decan2019empirical} broadened the lens across seven packaging ecosystems, characterising backward-incompatible updates and dependency staleness as a shared, structural hazard of large dependency networks rather than an ecosystem-specific accident. Collectively this literature establishes that the version string is a weak, statistically unreliable proxy for the property developers actually care about: whether \emph{their} call sites still behave correctly.

A second strand attacks the detection problem directly. Tools such as APIDiff~\cite{brito2018apidiff} and the framework Maracas underpinning Ochoa et al.~\cite{ochoa2022breaking} compare two library versions at the level of types, methods, and fields, computing a delta of signature-level changes and, in the case of Maracas, the set of client locations each change impacts. Xavier et al.~\cite{xavier2017historical} performed a large-scale historical study of such syntactic breaking changes and their client impact, and Brito et al.~\cite{brito2018why} catalogued the reasons developers introduce them. These signature-diffing techniques are mature and precise, but they share a definitional ceiling: they detect changes to the \emph{shape} of an interface, not to its \emph{behaviour}. A method whose signature is untouched but whose contract has silently shifted---an off-by-one boundary, a newly thrown exception, a changed null-handling policy---is invisible to them.

This ceiling motivates the literature on \emph{behavioural} backward incompatibility, which is the closest prior work to the present article. Mostafa et al.~\cite{mostafa2017behavioral} manually studied 296 behavioural backward incompatibilities across popular Java libraries, showing that changes with untouched signatures are both common and disproportionately damaging precisely because no diffing tool flags them. Foo et al.~\cite{foo2018efficient} proposed efficient static checking of library updates to surface such hazards without full re-execution, and Jezek and Dietrich~\cite{jezek2017magic} demonstrated experimentally that detection rates vary substantially across existing tools, with many real incompatibilities escaping every one. Most recently, Jayasuriya et al.~\cite{jayasuriya2024understanding} quantified the phenomenon at scale, finding that 11.58\% of dependency updates carry client-impacting breaking changes and that 2.30\% of updates introduce behavioural breaks detectable only when client tests are executed. The recurring methodological move in this strand is to recover behaviour dynamically---through the client's own test suite, regression runs, or learned invariants in the style of Daikon~\cite{ernst2007daikon}---which inherits the fundamental limitation of testing: it observes sampled inputs, never the whole input space, and so can confirm the presence of an incompatibility but never its absence. A complementary, recent direction uses large language models to \emph{repair} breaking dependency updates~\cite{islam2025byam} once detected, but presupposes that detection has already occurred and offers no machine-checkable guarantee that the repair preserves behaviour.

The decisive gap across all three strands is that compatibility is treated as an \emph{ex post}, sampled, and largely syntactic property: a version string is consulted, a signature diff is computed, or a test suite is re-run after the fact. None grounds compatibility in a machine-checkable behavioural contract that covers all inputs, and none answers the question a developer actually poses at upgrade time---``which of my call sites does this change break, and why?''---with a proof obligation rather than a failed test. The pipeline proposed here closes this gap by making the formal contract the unit of compatibility. Because contracts compose, a callee-spec change mechanically recomputes each caller's proof obligations and flags the broken call sites \emph{without running them}, turning behavioural-compatibility analysis into a static, exhaustive, counterexample-bearing operation. This builds directly on the compositional refinement established in our prior work~\cite{edmonds_article3}, in which a second weakest-precondition pass recovers preconditions a single pass misses, and on its integration into continuous-integration workflows~\cite{edmonds_article4}. In this framing an internal refactor and a third-party version bump are the \emph{same} operation---blast-radius analysis over a contract graph---and the foundational notion of an interface as a behavioural contract~\cite{meyer1992dbc} is finally enforced rather than merely documented.
